% =============================================================================
% ANEXO Q — Estructura de Repositorios
% =============================================================================
\chapter{Estructura de los Repositorios}
\label{cap:anexo_repos}

Árbol de directorios real de ambos repositorios, base de la organización \textit{vertical slice}
(backend) y \textit{feature-first} (frontend). El código fuente se publica en GitHub:
\url{https://github.com/dpardo23/EthosBack.git} (backend) y
\url{https://github.com/dpardo23/EthosFront.git} (frontend).\par

\section{Backend (\texttt{EthosBack})}
\label{sec:anexo_repo_be}
\begin{lstlisting}[caption={Estructura del repositorio backend.}, label={lst:anexo_be_tree}]
EthosBack/
|-- build.sh                  # Empaquetado monolitico (FE + BE -> JAR)
|-- pom.xml                   # Dependencias y plugins Maven
|-- Dockerfile / docker-compose.yml
`-- src/
    |-- main/
    |   |-- java/com/ethoshub/
    |   |   |-- EthosHubApplication.java
    |   |   |-- auth/ security/ profile/ project/ experience/
    |   |   |-- education/ skills/ portfolio/ cvstudio/
    |   |   |-- connections/ chat/ notifications/ preferences/
    |   |   |-- recruiter/ admin/
    |   |   |-- config/        # Seguridad, CORS, filtros
    |   |   `-- shared/        # ApiResponse, excepciones, util
    |   `-- resources/
    |       |-- application.yml / application-prod.yml
    |       |-- db/migration/V*.sql
    |       `-- static/        # dist del frontend (generado)
    `-- test/java/             # JUnit 5 + Testcontainers
\end{lstlisting}

\clearpage
\section{Frontend (\texttt{EthosFront})}
\label{sec:anexo_repo_fe}
\begin{lstlisting}[caption={Estructura del repositorio frontend.}, label={lst:anexo_fe_tree}]
EthosFront/
|-- index.html
|-- vite.config.ts            # Vite + proxy /api,/auth,/v1 -> :8080 + Vitest
|-- tailwind.config.js / tsconfig.json / eslint.config.js
`-- src/
    |-- main.tsx
    |-- app/                  # providers, layouts, router
    |-- pages/{auth,dashboard,recruiter,admin,public}/
    |-- features/{portfolio,projects,recruiter}/
    |-- components/{ui,auth,brand,preferences}/
    |-- shared/{api,services,types,lib,motion,landing,ui}/
    |-- store/                # Stores Zustand por dominio
    |-- hooks/
    `-- i18n/locales/{es,en,pt}/
\end{lstlisting}

\begin{NotaTecnica}
Ambos repositorios mantienen un grafo de conocimiento en \comando{graphify-out/}, consultable con
\comando{graphify query "<pregunta>"}, que documenta de forma navegable las relaciones del código.
\end{NotaTecnica}

\clearpage

% =============================================================================
\chapter{Convenciones de Código y Control de Versiones}
\label{cap:anexo_convenciones}

Convenciones transversales del proyecto, fijadas en el Sprint~0 y mantenidas en todos los
incrementos.\par

\section{Control de Versiones (GitFlow)}
\label{sec:anexo_conv_git}
\begin{itemize}[noitemsep]
    \item Ramas \comando{main} (producción), \comando{develop} (integración) y
    \comando{feature/*} (trabajo aislado).
    \item \textit{Tags} semánticos por sprint: \comando{release/v0.1.0} \dots \comando{release/v2.0.0}.
    \item Cada \textit{pull request} actualiza la documentación asociada (\textit{Documentation as Code}).
\end{itemize}

\section{Convenciones de Código}
\label{sec:anexo_conv_code}
\renewcommand{\arraystretch}{1.3}
\fontsize{9}{10.8}\selectfont\sffamily
\noindent
\begin{tabularx}{\textwidth}{|l|X|}
\hline
\rowcolor{colorPrimario}
\textcolor{white}{\textbf{Ámbito}} & \textcolor{white}{\textbf{Convención}} \\ \hline
Backend & Organización \textit{vertical slice}: \comando{controller+service+repository+dto}. \\ \hline
\rowcolor{grisTecnico}
DTO & \comando{record} validados con Jakarta Bean Validation. \\ \hline
Respuestas & Envoltorio común \comando{ApiResponse<T>}. \\ \hline
\rowcolor{grisTecnico}
Serialización & Jackson en UTC; \comando{default-property-inclusion: NON\_NULL}. \\ \hline
Frontend & TypeScript \comando{strict}; alias \comando{@/*} $\to$ \comando{src/*}. \\ \hline
\rowcolor{grisTecnico}
Lint & ESLint 10 (\textit{flat config}). \\ \hline
Estado & \textit{Stores} Zustand por dominio; \comando{resetAllStores} en logout. \\ \hline
\end{tabularx}\normalfont\normalsize

\clearpage
